
\documentclass{article}
\usepackage{colortbl}
\usepackage{makecell}
\usepackage{multirow}
\usepackage{supertabular}

\begin{document}

\newcounter{utterance}

\centering \large Interaction Transcript for game `wordle', experiment `high\_frequency\_words\_no\_clue\_no\_critic', episode 3 with gemma{-}4{-}26b{-}a4b{-}it.
\vspace{24pt}

{ \footnotesize  \setcounter{utterance}{1}
\setlength{\tabcolsep}{0pt}
\begin{supertabular}{c@{$\;$}|p{.15\linewidth}@{}p{.15\linewidth}p{.15\linewidth}p{.15\linewidth}p{.15\linewidth}p{.15\linewidth}}
    \# & \multicolumn{2}{c}{Player} && \multicolumn{2}{c}{Game Master} \\
    \hline

    \theutterance \stepcounter{utterance}  
    & & & \multicolumn{4}{p{0.6\linewidth}}{
        \cellcolor[rgb]{0.9,0.9,0.9}{
            \makecell[{{p{\linewidth}}}]{
                \texttt{\tiny{[P1$\langle$GM]}}
                \texttt{Eres un mago del lenguaje al que le gusta adivinar palabras usando las reglas dadas.} \\
\\ 
\texttt{¡Bienvenido a Wordle! Tienes seis intentos para adivinar la palabra objetivo, una palabra válida en español de cinco letras minúsculas. Por favor, usa las etiquetas "explicación:" e "intento:" para proporcionar una explicación concisa para cada intento.} \\
\\ 
\texttt{Por ejemplo, si tu intento es "parte", tu respuesta debería ser} \\
\texttt{explicación: esta es una palabra común en español de cinco letras, y estoy comenzando mi intento con esta palabra.} \\
\texttt{intento: parte} \\
\\ 
\texttt{Después de cada intento, tu respuesta será validada y recibirás retroalimentación indicando qué letras son correctas (verde), qué letras son correctas pero están en la posición incorrecta (amarillo) y qué letras son incorrectas (rojo). Esta retroalimentación puede ser útil para determinar qué letras incluir o excluir en tu siguiente intento.} \\
\\ 
\texttt{Por ejemplo, la retroalimentación para "parte" podría ser:} \\
\texttt{retroalimentación\_intento: p<amarillo> a<amarillo> r<verde> t<amarillo> e<rojo>} \\
\\ 
\texttt{La explicación debe contener detalles sobre cómo se usa el retroalimentación\_intento para llegar a un nuevo intento.} \\
\\ 
\texttt{No generes ningún otro texto.} \\
\\ 
\texttt{Comencemos con tu primer intento.} \\
            }
        }
    }
     \\ \\

    \theutterance \stepcounter{utterance}  
    & \multicolumn{4}{p{0.6\linewidth}}{
        \cellcolor[rgb]{0.561,0.737,0.561}{
            \makecell[{{p{\linewidth}}}]{
                \texttt{\tiny{[P1$\rangle$GM]}}
                \texttt{explicación: para comenzar mi búsqueda, elegiré una palabra con vocales comunes y consonantes frecuentes en español para maximizar la información obtenida.} \\
\texttt{intento: audio} \\
            }
        }
    }
    & & \\ \\

    \theutterance \stepcounter{utterance}  
    & & & \multicolumn{2}{p{0.3\linewidth}}{
        \cellcolor[rgb]{0.9,0.9,0.9}{
            \makecell[{{p{\linewidth}}}]{
                \texttt{\tiny{[GM$|$GM]}}
                \texttt{attempts = 1} \\
\texttt{target\_word = hecho} \\
\texttt{guess = audio} \\
\texttt{guess\_feedback = a<red> u<red> d<red> i<red> o<green>} \\
            }
        }
    }
    & & \\ \\

    \theutterance \stepcounter{utterance}  
    & & & \multicolumn{4}{p{0.6\linewidth}}{
        \cellcolor[rgb]{0.9,0.9,0.9}{
            \makecell[{{p{\linewidth}}}]{
                \texttt{\tiny{[P1$\langle$GM]}}
                \texttt{retroalimentación\_intento: a<rojo> u<rojo> d<rojo> i<rojo> o<verde>} \\
\\ 
\texttt{Proporciona tu respuesta solo en este formato.} \\
\texttt{explicación: detalles} \\
\texttt{intento: palabra} \\
            }
        }
    }
     \\ \\

    \theutterance \stepcounter{utterance}  
    & \multicolumn{4}{p{0.6\linewidth}}{
        \cellcolor[rgb]{0.561,0.737,0.561}{
            \makecell[{{p{\linewidth}}}]{
                \texttt{\tiny{[P1$\rangle$GM]}}
                \texttt{explicación: la única letra correcta es la 'o' en la quinta posición. las letras 'a', 'u', 'd' e 'i' han sido descartadas. buscaré una palabra que termine en 'o' con otras consonantes y vocales comunes como 'e' o 'l'.} \\
\texttt{intento: menos} \\
            }
        }
    }
    & & \\ \\

    \theutterance \stepcounter{utterance}  
    & & & \multicolumn{2}{p{0.3\linewidth}}{
        \cellcolor[rgb]{0.9,0.9,0.9}{
            \makecell[{{p{\linewidth}}}]{
                \texttt{\tiny{[GM$|$GM]}}
                \texttt{attempts = 2} \\
\texttt{target\_word = hecho} \\
\texttt{guess = menos} \\
\texttt{guess\_feedback = m<red> e<green> n<red> o<yellow> s<red>} \\
            }
        }
    }
    & & \\ \\

    \theutterance \stepcounter{utterance}  
    & & & \multicolumn{4}{p{0.6\linewidth}}{
        \cellcolor[rgb]{0.9,0.9,0.9}{
            \makecell[{{p{\linewidth}}}]{
                \texttt{\tiny{[P1$\langle$GM]}}
                \texttt{retroalimentación\_intento: m<rojo> e<verde> n<rojo> o<amarillo> s<rojo>} \\
\\ 
\texttt{Proporciona tu respuesta solo en este formato.} \\
\texttt{explicación: detalles} \\
\texttt{intento: palabra} \\
            }
        }
    }
     \\ \\

    \theutterance \stepcounter{utterance}  
    & \multicolumn{4}{p{0.6\linewidth}}{
        \cellcolor[rgb]{0.561,0.737,0.561}{
            \makecell[{{p{\linewidth}}}]{
                \texttt{\tiny{[P1$\rangle$GM]}}
                \texttt{explicación: la letra 'e' está en la segunda posición y la 'o' debe estar en la palabra pero no en la quinta posición. las letras 'm', 'n' y 's' han sido descartadas. probaré con una palabra que use la 'o' en otra posición y nuevas consonantes como 'p', 'l' o 'r'.} \\
\texttt{intento: pelo} \\
            }
        }
    }
    & & \\ \\

    \theutterance \stepcounter{utterance}  
    & & & \multicolumn{2}{p{0.3\linewidth}}{
        \cellcolor[rgb]{0.9,0.9,0.9}{
            \makecell[{{p{\linewidth}}}]{
                \texttt{\tiny{[GM$|$GM]}}
                \texttt{Error: The length of the guessed word is not 5.} \\
            }
        }
    }
    & & \\ \\

    \theutterance \stepcounter{utterance}  
    & & & \multicolumn{2}{p{0.3\linewidth}}{
        \cellcolor[rgb]{0.9,0.9,0.9}{
            \makecell[{{p{\linewidth}}}]{
                \texttt{\tiny{[GM$|$GM]}}
                \texttt{game\_result = ABORT} \\
            }
        }
    }
    & & \\ \\

\end{supertabular}
}

\end{document}
